\section{ПРОГРАММНАЯ РЕАЛИЗАЦИЯ}

\subsection{Выбор и обоснование технологического стека}

На основе спроектированной в предыдущей главе архитектуры была выполнена программная реализация системы. Выбор технологического стека определялся требованиями к производительности, поддерживаемости и удобству разработки. Для серверной части был выбран язык Python версии~3.11 в сочетании с фреймворком FastAPI~\cite{fastapi_docs}. Python обеспечивает высокую скорость разработки благодаря развитой экосистеме библиотек для работы с базами данных, статистического анализа и визуализации. FastAPI предоставляет современный асинхронный веб-фреймворк с автоматической генерацией документации API через OpenAPI, встроенной валидацией данных через Pydantic версии~2 и высокой производительностью, сопоставимой с фреймворками на Node.js. Управление конфигурацией осуществляется через библиотеку pydantic-settings, загружающую параметры из переменных окружения и файла \texttt{.env}.

Для асинхронного выполнения запросов к СУБД выбраны драйверы asyncpg для PostgreSQL и aiomysql для MySQL и MariaDB. Оба драйвера обеспечивают неблокирующий ввод-вывод, что позволяет эффективно обрабатывать множество параллельных соединений в рамках одного потока операционной системы. В качестве ORM-слоя используется SQLAlchemy версии 2.0 с поддержкой асинхронных сессий, что обеспечивает типобезопасную работу с базой данных истории. Для управления миграциями схемы базы данных применяется инструмент Alembic, позволяющий описывать изменения схемы в виде версионируемых Python-скриптов и применять их автоматически при старте приложения.

Для клиентской части выбран фреймворк Next.js версии~16 на базе React~19 с использованием TypeScript. Next.js обеспечивает клиентскую маршрутизацию через App Router и встроенную поддержку серверных компонентов. TypeScript гарантирует типобезопасность на всех уровнях клиентского приложения, снижая количество ошибок времени выполнения. Для управления состоянием приложения выбрана библиотека Zustand, отличающаяся минимальным объёмом кода и отсутствием необходимости в оборачивании приложения в провайдеры. Для визуализации данных используются библиотеки Recharts и Plotly: Recharts предоставляет декларативный API для построения интерактивных графиков на базе SVG, Plotly применяется для диаграмм размаха и расширенных статистических визуализаций. Стилизация реализована через Tailwind CSS версии~4 с CSS-переменными как дизайн-токенами; компоненты интерфейса построены на базе shadcn/ui поверх Radix UI, что обеспечивает доступность и корректную работу с клавиатурой.

Для развёртывания всех компонентов системы используется Docker Compose; состав сервисов, порты и особенности сборки образов рассматриваются в разделе о контейнеризации настоящей главы.

\subsection{Реализация серверной части}

Серверная часть системы реализована в виде модульного приложения FastAPI. Точка входа находится в файле \texttt{main.py}, где инициализируется экземпляр приложения, регистрируются маршруты и создаётся WebSocket-эндпоинт.

Инициализация компонентов системы выполняется через модуль \texttt{initialize.py}, который создаёт экземпляры пяти репозиториев: \texttt{TestRepository}, \texttt{ConnectionRepository}, \texttt{ProfileRepository}, \texttt{ScenarioBundleRepository} и \texttt{DatabaseGroupRepository}. Каждый репозиторий принимает URL базы данных истории и создаёт собственный пул асинхронных соединений через SQLAlchemy. Настройки стратегии восстановления баз данных \texttt{settings.restore.*} загружаются из модуля конфигурации \texttt{core/config.py} и изменяются через \texttt{settings\_routes} без отдельного репозитория. Такой подход реализует паттерн внедрения зависимостей и обеспечивает централизованное управление жизненным циклом компонентов. При недоступности базы данных инициализация не прерывает запуск приложения~--- репозитории устанавливаются в \texttt{None}, а соответствующая функциональность корректно деградирует.

Запуск и остановка приложения управляются через менеджер контекста lifespan. При запуске выполняются два последовательных шага:

\begin{enumerate}[label=\arabic*)]
    \item автоматически применяются миграции Alembic командой \texttt{upgrade head}, что гарантирует актуальность схемы базы данных вне зависимости от версии развёртывания;
    \item при успешной инициализации всех трёх ключевых репозиториев: \texttt{connection\_repository}, \texttt{profile\_repository} и \texttt{scenario\_bundle\_repository}~--- запускается \texttt{LogicalScenarioBootstrap}: он создаёт шесть встроенных шаблонов сценариев и профили схем, а при наличии зарегистрированных подключений с определёнными профилями~--- генерирует SQL-наборы для пяти шаблонов с автогенерацией и встроенные OLTP-наборы для поддерживаемых профилей.
\end{enumerate}

Первый шаг является fail-fast: при возникновении ошибки Alembic запуск приложения немедленно прерывается, что исключает работу системы с устаревшей схемой базы данных. Bootstrap на втором шаге выполняется условно~--- только при наличии рабочего подключения к базе данных проекта; если база данных недоступна, bootstrap пропускается без прерывания запуска. Повторный запуск обоих шагов не создаёт дублирующих записей~--- процедура полностью идемпотентна.

\subsubsection{Модуль генерации нагрузки}

Модуль генерации нагрузки реализован в классе \texttt{LoadTester} файла \texttt{load\_tester/tester.py}. Основным методом запуска является \texttt{run\_resolved\_scenario\_test\_suite}, принимающий разрешённый набор сценария нагрузки и список идентификаторов подключений. В зависимости от типа набора метод выполняет либо одиночные SQL-запросы, либо транзакционные последовательности операций, где каждый виртуальный пользователь атомарно выполняет группу инструкций изменения данных. Единая структура каждого прогона: подготовка базы данных, двухфазный прогрев, измерительная фаза, сбор метрик, восстановление состояния, верификация. Блок \texttt{try/finally} гарантирует выполнение восстановления даже при возникновении ошибок в ходе тестирования.

Базовый метод \texttt{execute\_query} выполняет отдельный SQL-запрос к указанной базе данных, замеряя время отклика через \texttt{time.perf\_counter()}. На рисунке~\ref{lst:execute_query} приведён фрагмент реализации:

\begin{figure}
    \begin{minipage}{\linewidth}
\begin{lstlisting}
async def execute_query(self, db_key, query, query_id):
    """ Выполнение одного запроса с измерением времени """
    start_time = time.perf_counter()
    error = None
    rows_count = 0
    db_type = self.db_connection.get_dbms_type(db_key)

    try:
        engine = await self.db_connection.get_engine_async(db_key)
        async with engine.connect() as conn:
            result = await conn.execute(text(query))
            rows_count = len(result.fetchall()) if result.returns_rows else 0
            await conn.commit()
    except Exception as e:
        error = str(e)

    end_time = time.perf_counter()
    execution_time = (end_time - start_time) * 1000

    return {
        'query_id': query_id,
        'execution_time_ms': execution_time,
        'rows_count': rows_count,
        'error': error,
        'timestamp': datetime.now(timezone.utc).isoformat()
    }
\end{lstlisting}
    \end{minipage}
    \caption{Реализация метода \texttt{execute\_query} класса LoadTester}
    \label{lst:execute_query}
\end{figure}

Использование \texttt{time.perf\_counter()} обеспечивает высокую точность замеров; асинхронное подключение через SQLAlchemy engine позволяет выполнять запросы без блокировки потока; результат включает не только время выполнения, но и количество обработанных строк и текст ошибки при её возникновении.

Каждый виртуальный пользователь реализован как асинхронная задача \texttt{asyncio.Task}. На каждой итерации выполняется одиночный запрос или транзакционная последовательность операций в соответствии с \texttt{workload\_mode} разрешённого набора сценария; результаты накапливаются в буфере. При достижении интервала в одну секунду вычисляются агрегированные метрики окна: среднее время отклика, успешная и полная пропускная способность, количество успешных и ошибочных операций; итог прогона дополняется перцентилями p50, p95, p99, минимумом, максимумом и стандартным отклонением за измерительную фазу. Данные передаются через callback-механизм для трансляции в WebSocket и сохранения в базу данных истории.

Для сбора системных метрик используется метод \texttt{get\_system\_metrics}, который через библиотеку \texttt{psutil} получает загрузку центрального процессора (проценты), использование оперативной памяти (проценты и мегабайты), дисковые операции в секунду, скорости чтения и записи (МиБ/с), а также входящий и исходящий сетевой трафик (МиБ/с). Метод \texttt{get\_dbms\_metrics} через слой диалектов в каталоге \texttt{database/dialects/} возвращает унифицированный набор внутренних показателей: коэффициент попаданий в кэш, размер буферного пула (мегабайты и метка параметра), число ожиданий блокировок и взаимных блокировок, объёмы таблиц. Для PostgreSQL используются представления \texttt{pg\_stat\_database}~\cite{postgresql_docs} и \texttt{pg\_stat\_bgwriter}, для MySQL и MariaDB~--- \texttt{performance\_schema.data\_lock\_waits} и \texttt{SHOW GLOBAL STATUS}. Слой диалектов изолирует СУБД-специфичную логику и позволяет расширять поддержку СУБД без изменения основного цикла сбора метрик.

\subsubsection{Модуль генерации сценариев}

Модуль генерации сценариев реализован в классе \texttt{ScenarioGenerator} файла \texttt{database/scenario\_generator.py}. Генератор принимает на вход подключение к целевой СУБД и тип шаблона сценария, анализирует метаданные схемы через \texttt{SchemaAnalyzer} и формирует структуру \texttt{ScenarioBundle} с массивами запросов и индексов.

Для пяти встроенных шаблонов с автогенерацией, кроме OLTP, \texttt{ScenarioGenerator} строит параметризованные SQL-шаблоны: шаблон только чтения включает выборки по первичному ключу и JOIN-запросы по внешним ключам, только записи~--- операции INSERT и UPDATE на изменяемых таблицах, смешанные шаблоны~--- комбинации SELECT и UPDATE с различными весами, аналитические запросы~--- агрегации, JOIN и диапазонные чтения по метаданным схемы. Шаблон OLTP не проходит через генератор: для профилей Sakila и Brazilian E-Commerce наборы транзакций задаются встроенными конфигурациями. Параметры запросов типа \texttt{random\_from\_table} ссылаются на реальные столбцы схемы и кэшируются по ключу \texttt{db\_key:table:column}. Класс \texttt{ScenarioBundleResolver} разрешает конкретный набор сценария нагрузки при запуске теста: по идентификатору набора или по паре <<профиль + шаблон>> он загружает запросы, параметры и индексы, формируя неизменяемый снимок конфигурации, который сохраняется вместе с прогоном для воспроизводимости.

\subsubsection{Callback-механизм потоковой передачи данных}

Ниже детализирован callback-механизм потоковой передачи метрик, упомянутый при описании \texttt{LoadTester}. Для связи с WebSocket и базой данных истории используется паттерн Observer. Класс \texttt{LoadTester} поддерживает два независимых callback: callback метрик, вызываемый на каждой итерации сбора агрегированных показателей, и callback статуса операций сохранения/восстановления состояния базы данных. Такое разделение позволяет клиенту независимо отслеживать ход теста и состояние данных.

Класс \texttt{TestStreamingCallback} файла \texttt{websocket\_manager.py} выступает в роли наблюдателя: он получает метрики от \texttt{LoadTester} и транслирует их подписчикам. Метод формирует объект \texttt{TestMetricsUpdate}, содержащий все метрики производительности, системные показатели и прогресс выполнения. Данные передаются через \texttt{ConnectionManager} для рассылки всем WebSocket-подписчикам теста и одновременно сохраняются в базе данных истории как точка временного ряда. Для оптимизации записи сырых выборок метрик используется буферизация: накопление происходит до достижения размера буфера в 10 записей, после чего выполняется пакетная вставка через метод \texttt{add\_metric\_sample\_batch}.

\subsubsection{Управление WebSocket-соединениями}

Управление WebSocket-соединениями реализовано в классе \texttt{ConnectionManager} файла \texttt{websocket\_manager.py}. Класс хранит активные соединения в словаре \texttt{active\_connections}, ключом которого является идентификатор теста, а значением~--- список WebSocket-объектов. Метод \texttt{connect} принимает новое соединение и добавляет его в список, метод \texttt{disconnect} удаляет соединение и очищает пустые списки. Метод \texttt{broadcast\_to\_test} рассылает JSON-сообщение всем подписчикам конкретного теста. WebSocket-эндпоинт зарегистрирован по пути \texttt{/ws/test/\{test\_id\}} и обрабатывает входящие ping-сообщения от клиента для поддержания соединения.

\subsubsection{Управление состоянием базы данных}

Оркестратор управления состоянием баз данных реализован в классе \texttt{DatabaseStateManager} файла \texttt{database/state\_manager.py}. Класс управляет полным жизненным циклом операций: анализ запросов сценария, сохранение состояния базы данных, выполнение теста, восстановление исходных данных и верификация целостности.

Метод \texttt{prepare\_for\_test} сначала передаёт SQL-запросы сценария в модуль \texttt{QueryAnalyzer}, который определяет наличие write-операций INSERT, UPDATE и DELETE. Если сценарий содержит только SELECT-запросы, подготовка пропускается. При обнаружении write-операций определяются затронутые таблицы и полный список всех таблиц базы данных. Затем модуль \texttt{StateVerifier} фиксирует объект \texttt{StateFingerprint}~--- количество строк в каждой таблице, значения последовательностей и, при включённой настройке \texttt{verify\_checksums}, контрольные суммы данных,~--- который впоследствии используется для верификации восстановления. После этого через настроенную стратегию сохраняется состояние всех таблиц базы данных: не только затронутых write-операциями, но и всех остальных, что гарантирует полный откат к исходному состоянию.

Нативная стратегия \texttt{NativeDumpStrategy} является основной: она вызывает CLI-утилиты конкретной СУБД~--- \texttt{pg\_dump} для создания дампа и \texttt{pg\_restore}/\texttt{psql} для восстановления у PostgreSQL; \texttt{mysqldump}/\texttt{mysql} для MySQL; \texttt{mariadb-dump}/\texttt{mariadb} для MariaDB,~--- сохраняя результат в виде файла дампа в директории \texttt{snapshots}. Для MariaDB используется самостоятельный клиент \texttt{mariadb-dump} с приоритетом перед \texttt{mysqldump}: это обеспечивает корректную обработку DEFINER-клауз и специфических расширений синтаксиса. При восстановлении через нативную стратегию для MySQL/MariaDB дамп предварительно очищается от DEFINER-клауз, а для PostgreSQL выполняется TRUNCATE целевых таблиц перед data-only восстановлением через \texttt{pg\_restore}/\texttt{psql}. Если CLI-утилиты не обнаружены в PATH, \texttt{NativeDumpStrategy} автоматически переключается на \texttt{SqlBackupStrategy}.

SQL-стратегия \texttt{SqlBackupStrategy} создаёт внутрибазовые копии таблиц через запросы вида \texttt{CREATE~TABLE~\ldots~AS~SELECT~*} с фиксированным префиксом \texttt{\_loadtest\_backup\_} в той же схеме. При восстановлении таблицы очищаются через TRUNCATE и заполняются данными из копий через INSERT; порядок восстановления определяется топологической сортировкой по FK-зависимостям с временным отключением FK-ограничений. Данная стратегия не требует внешних утилит и работает через SQLAlchemy с любой поддерживаемой СУБД, однако временно увеличивает объём данных в тестируемой базе.

Метод \texttt{restore\_after\_test} выполняет восстановление под блокировкой \texttt{asyncio.Lock}, специфичной для типа СУБД, что исключает конфликты при параллельном тестировании. При включённой настройке \texttt{verify\_after\_restore} модуль \texttt{StateVerifier} строит новый \texttt{StateFingerprint} и сравнивает его с исходным: при полном совпадении числа строк и контрольных сумм восстановление считается успешным, при несовпадении формируется отчёт о расхождениях. По завершении операции внутрибазовые копии или временные файлы дампов удаляются.

\subsubsection{Разрешение наборов сценариев нагрузки и запуск теста}

Фоновая задача \texttt{run\_test\_with\_streaming} создаёт экземпляр \texttt{LoadTester}, устанавливает callback для потоковой передачи метрик и callback статуса операций сохранения/восстановления состояния, разрешает набор сценария нагрузки через \texttt{ScenarioBundleResolver} и передаёт его в \texttt{LoadTester} вместе со списком идентификаторов подключений. Перед началом теста \texttt{IndexManager} создаёт индексы, описанные в \texttt{ScenarioBundleIndex}; по завершении теста индексы удаляются. По завершении результаты собираются, сохраняются в базе данных, а статус обновляется на \texttt{completed} или \texttt{failed}. При передаче параметра \texttt{database\_group\_id} тестовый прогон автоматически привязывается к соответствующей группе баз данных.

\subsubsection{Модуль сравнительного анализа}

Модуль сравнительного анализа реализован в классе \texttt{ComparisonService} файла \texttt{comparison/service.py}. Конструктор класса принимает три репозитория: \texttt{repository} для загрузки данных тестов, \texttt{scenario\_bundle\_repository} для получения информации о наборах сценариев нагрузки и \texttt{connection\_repository} для получения данных о подключениях. Основной метод \texttt{analyze} принимает объект \texttt{ComparisonRequest} с полями \texttt{analysis\_mode}, \texttt{test\_ids}, необязательным \texttt{baseline\_id} и конфигурацией отчёта.

При обнаружении потенциальных проблем сопоставимости~--- разных наборов сценариев нагрузки, неравных временных периодов, смешанных типов СУБД~--- система формирует предупреждения, включаемые в итоговый отчёт.

Метод \texttt{analyze} маршрутизирует запрос по полю \texttt{analysis\_mode}, задаваемому клиентом на странице истории тестов. В режиме \texttt{per\_test} передаётся ровно один идентификатор прогона: загружаются его конфигурация, результаты и сырые выборки метрик по каждому подключению с ключом \texttt{connection\_key}; для всех пар СУБД внутри прогона вычисляется описательная статистика задержки и пропускной способности, применяются статистические критерии с FDR-поправкой и формируется ранжирование; отчёт генерирует \texttt{PerTestReportGenerator}. В режиме \texttt{series} передаётся от двух до пяти идентификаторов прогонов с различными уровнями нагрузки: для каждой СУБД строятся траектории метрик, оцениваются показатели масштабируемости \texttt{stability\_index}, \texttt{elasticity} и точка насыщения, тренды проверяются критериями Спирмена и Манна--Кендалла; базовый прогон задаётся полем \texttt{baseline\_id}; отчёт генерирует \texttt{SeriesReportGenerator}. Оба генератора реализованы в модуле \texttt{analysis/report\_generator.py}.

Сбор сырых выборок метрик реализован с трёхуровневым механизмом fallback: при наличии записей \texttt{MetricSample} используются они; при их отсутствии~--- данные из \texttt{TimeSeries}; в крайнем случае~--- агрегированные метрики из \texttt{TestResult} с соответствующим предупреждением о сниженной надёжности анализа. Это обеспечивает работоспособность сравнения даже для тестов, выполненных до добавления модуля сбора сырых метрик.

Отчёт в режиме \texttt{per\_test} включает краткое резюме с лидером по пропускной способности и СУБД с наименьшей средней задержкой, попарные сравнения с интерпретацией и ранжирование подключений. Отчёт в режиме \texttt{series} дополняется блоком \texttt{parameter\_impacts} с анализом влияния различий конфигурации между прогонами; в обоих режимах автоматически выявляются паттерны нестабильности при высоком CV, деградации при масштабировании и высоких хвостовых задержек при соотношении p99 к медиане.

\subsubsection{Реализация статистических функций}

Модуль \texttt{comparison/statistics.py} содержит математический аппарат для статистического сравнения выборок. Функция \texttt{calculate\_descriptive\_stats} вычисляет описательную статистику через \texttt{numpy}: количество, среднее, медиану, стандартное отклонение, минимум, максимум, процентили p50, p95, p99, коэффициент вариации CV, а также межквартильный размах IQR. Функция \texttt{check\_normality} проверяет нормальность распределения через критерий Шапиро--Уилка с ограничением выборки в 5000 элементов.

Функция \texttt{compare\_two\_samples} реализует двухэтапный алгоритм выбора критерия. На первом этапе проверяется размер выборки: если количество элементов менее~10, статистический тест не выполняется и выводится предупреждение о недостаточном объёме данных; если размер выборки превышает 5000~элементов, из неё извлекается случайная подвыборка для снижения вычислительной нагрузки. Затем к полученной выборке применяется критерий Шапиро--Уилка: если p-значение превышает 0,05, нулевая гипотеза о нормальности не отвергается. На втором этапе по результату проверки нормальности выбирается критерий сравнения: при нормальном распределении применяется критерий Уэлча, не требующий равенства дисперсий; при ненормальном~--- критерий Манна--Уитни, являющийся непараметрическим аналогом t-критерия. Если p-значение выбранного критерия оказывается меньше~0,05, различие признаётся статистически значимым и формируется текстовая интерпретация с указанием процентного отклонения. При сравнении нескольких пар применяется FDR-поправка Бенджамини--Хохберга для контроля уровня ложных открытий.

\subsubsection{Реализация REST API}

REST API системы разделено на восемь модулей маршрутов, каждый из которых отвечает за свою область функциональности. Основные эндпоинты представлены в таблице~\ref{tab:api_endpoints}.

\begin{table}
    \centering
    \caption{Основные эндпоинты REST API}
    \label{tab:api_endpoints}
    \begin{tabularx}{\textwidth}{|>{\RaggedRight}X|>{\RaggedRight}X|>{\RaggedRight}X|}
        \hline
        Модуль & Метод и путь & Назначение \\
        \hline
        test & POST /test/async & Асинхронный запуск теста \\
        \hline
        test & GET /test/async/\{id\} & Получение статуса теста \\
        \hline
        test & GET /test/async/\{id\}/results & Получение результатов \\
        \hline
        history & GET /history/tests & Список тестов с пагинацией \\
        \hline
        history & GET /history/tests/\{id\} & Детали конкретного теста \\
        \hline
        history & DELETE /history/tests/\{id\} & Удаление теста \\
        \hline
        connection & GET /api/connections & Список подключений \\
        \hline
        connection & POST /api/connections & Создание подключения \\
        \hline
        connection & POST /api/connections/test & Тестирование подключения \\
        \hline
        database\_state & POST /.../\{id\}/backup & Сохранение состояния базы данных \\
        \hline
        database\_state & POST /.../\{id\}/restore & Восстановление состояния базы данных \\
        \hline
        database\_state & GET /.../\{id\}/estimate & Оценка объёма данных базы \\
        \hline
        comparison & POST /api/comparison/analyze & Сравнительный анализ \\
        \hline
        schema-profiles & GET /api/schema-profiles & Список профилей схемы \\
        \hline
    \end{tabularx}
\end{table}

\begin{table}
    \centering
    \ContinuedFloat
    \caption*{Продолжение таблицы \ref{tab:api_endpoints}}
    \label{tab:api_endpoints_continue}
    \begin{tabularx}{\textwidth}{|>{\RaggedRight}X|>{\RaggedRight}X|>{\RaggedRight}X|}
        \hline
        schema-profiles & POST /api/schema-profiles & Создание профиля \\
        \hline
        schema-profiles & GET /.../\{id\}/bundles & наборы сценариев нагрузки профиля \\
        \hline
        schema-profiles & POST /.../\{id\}/bundles/generate & Генерация набора сценария нагрузки для профиля схемы \\
        \hline
        database-groups & GET /api/database-groups & Список групп баз данных \\
        \hline
        database-groups & POST /api/database-groups & Создание группы баз данных \\
        \hline
        database-groups & PUT /.../\{id\}/profile & Привязка профиля схемы к группе \\
        \hline
        database-groups & POST /.../\{id\}/bundles/generate & Генерация набора сценария нагрузки для профиля, привязанного к группе баз данных \\
        \hline
        database-groups & POST /.../\{id\}/connections/\allowbreak\{connection\_id\}/confirm-profile & Подтверждение профиля схемы для подключения в группе \\
        \hline
        database-groups & PUT /.../\{id\}/reference-connection & Обновление эталонного подключения \\
        \hline
        settings & GET /api/settings/restore & Настройки восстановления \\
        \hline
        settings & PUT /api/settings/restore & Обновление настроек восстановления \\
        \hline
    \end{tabularx}
\end{table}

Эндпоинты \texttt{POST /.../\{id\}/bundles/generate} в модулях \texttt{schema-profiles} и \texttt{database-groups} вызывают одну операцию генерации наборов сценариев нагрузки через \texttt{ScenarioGenerator}; различаются контекст запроса и проверки входных параметров.

Эндпоинт \texttt{POST /test/async} принимает параметры нагрузки, создаёт запись прогона со статусом \texttt{pending}, запускает фоновую задачу \texttt{run\_test\_with\_streaming} через \texttt{BackgroundTasks} FastAPI и возвращает клиенту \texttt{test\_id} с \texttt{websocket\_url} для подписки на метрики. Настройки стратегии восстановления по пути \texttt{/api/settings/restore} изменяются в памяти процесса до рестарта бэкенда.

\subsubsection{Обработка ошибок}

Обработка ошибок в серверной части реализована на нескольких уровнях. На уровне API маршрутов все исключения перехватываются и преобразуются в HTTP-статусы: \texttt{ValueError}~--- в 400 Bad Request, неожиданные исключения~--- в 500 Internal Server Error. На уровне бизнес-логики ошибки логируются с указанием контекста и, где возможно, обрабатываются без прерывания выполнения теста. Блок \texttt{try/finally} в методах тестирования гарантирует выполнение восстановления базы данных даже при критических ошибках.

При возникновении ошибки выполнения SQL-запроса метод \texttt{execute\_query} не прерывает цикл итераций, а сохраняет текст ошибки в результате и продолжает выполнение следующих запросов. Это позволяет получить полную картину производительности, включая частоту ошибок, которая сама по себе является важной метрикой.

\subsubsection{Слой доступа к данным}

Слой доступа к данным реализован через паттерн Repository. Базовый класс \texttt{BaseRepository} обеспечивает асинхронное подключение к базе данных истории через SQLAlchemy с управлением сессиями через контекстный менеджер. Наследники реализуют специфичные операции для своих сущностей: \texttt{TestRepository} управляет тестовыми прогонами и временными рядами, \texttt{ConnectionRepository}~--- конфигурациями подключений с расшифровкой паролей, \texttt{ProfileRepository}~--- профилями схем и шаблонами сценариев, \texttt{ScenarioBundleRepository}~--- наборами сценариев нагрузки с их запросами и индексами, \texttt{DatabaseGroupRepository}~--- группами баз данных.

Метод \texttt{TestRepository.add\_metric\_sample\_batch} выполняет пакетное сохранение сырых выборок метрик. Код метода приведён на рисунке~\ref{lst:save_metric_samples}:

\begin{figure}
    \begin{minipage}{\linewidth}
\begin{lstlisting}
async def add_metric_sample_batch(
    self, test_run_id: str, samples: List[Dict]
):
    """ Пакетное сохранение сырых метрик """
    if not samples:
        return

    async with self.get_session() as session:
        records = [
            MetricSample(
                test_run_id=test_run_id,
                db_type=s['db_type'],
                connection_key=s['connection_key'],
                query_id=s.get('query_id'),
                sample_type=s.get('sample_type'),
                timestamp=s['timestamp'],
                latency_ms=s.get('latency_ms'),
                throughput=s.get('throughput'),
                is_error=s.get('is_error', False),
                error_message=s.get('error_message'),
            )
            for s in samples
        ]
        session.add_all(records)
        await session.commit()
\end{lstlisting}
    \end{minipage}
    \caption{Метод addMetricSampleBatch класса TestRepository}
    \label{lst:save_metric_samples}
\end{figure}

Пакетная вставка снижает количество обращений к базе данных при высокой интенсивности тестирования. Каждая запись содержит тип выборки, что позволяет гибко выбирать уровень детализации данных при сравнительном анализе.

\subsubsection{Управление подключениями к СУБД}

Класс \texttt{DatabaseConnection} файла \texttt{database/connection.py} управляет подключениями к тестируемым СУБД. Подключения хранятся в репозитории базы данных истории и загружаются лениво при первом асинхронном вызове.

Метод \texttt{\_resolve\_connection\_key} обеспечивает поддержку shorthand-ключей: если передан ключ \texttt{mysql} или \texttt{postgresql} и существует ровно одно активное подключение данного типа, оно используется автоматически. При наличии нескольких подключений одного типа требуется явное указание \texttt{connection\_id}. Кэширование конфигураций подключений в памяти исключает повторные обращения к базе данных при каждом запросе.

Для шифрования паролей подключений используется модуль \texttt{core/crypto.py}, реализующий симметричное шифрование в формате Fernet. Пароль шифруется перед сохранением в базе данных и расшифровывается при установлении подключения к тестируемой СУБД. Ключ шифрования задаётся через переменную окружения \texttt{ENCRYPTION\_KEY}.

\subsection{Реализация клиентской части}

Клиентская часть системы реализована как одностраничное приложение на базе Next.js~16 с клиентской маршрутизацией. Главная страница в файле \texttt{app/page.tsx} определяет компонент-маршрутизатор, который условно отображает один из восьми компонентов-страниц в зависимости от значения \texttt{currentPage} в Zustand store. Общая компоновка включает верхнюю панель с логотипом и навигацией, боковое меню с переключением страниц и основную область контента.

\subsubsection{Архитектура компонентов}

Страница \texttt{HomePage} является стартовой и предоставляет краткое описание системы с кнопками быстрого перехода к конфигурации теста или истории. \texttt{ConnectionsPage} обеспечивает полный CRUD для сохранённых подключений к СУБД: добавление, редактирование, удаление, проверка доступности. При создании подключения страница отображает результат автоопределения профиля схемы с коэффициентом уверенности.

Страница \texttt{ConfigPage} является центральным узлом конфигурации теста. Она содержит несколько карточек: выбор группы баз данных или отдельных подключений, выбор режима теста, выбор сценария или набор сценария нагрузки, настройка параметров нагрузки: числа виртуальных пользователей, числа итераций и времени прогрева, а также сводка конфигурации и кнопка запуска.

Страница \texttt{ScenariosPage} предоставляет браузер наборов сценариев нагрузки через компонент \texttt{logical-scenarios-browser}. В браузере отображаются группы баз данных, профили схем и связанные с ними наборы сценариев нагрузки с их SQL-запросами, параметрами и спецификацией индексов. Предусмотрены операции создания, редактирования, клонирования и генерации наборов сценариев нагрузки из эталонного подключения. Для каждой группы баз данных отображаются статусы \texttt{profile\_status} и \texttt{compatibility\_status}, что позволяет пользователю контролировать состояние привязки профиля и совместимость подключённых баз данных.

Страница \texttt{DashboardsPage} предоставляет дашборд реального времени с четырьмя вкладками: время отклика и пропускная способность, внутренние метрики СУБД, объёмы данных и системные метрики хоста. Каждая вкладка отображает секундные временны\'е ряды для всех подключений группы одновременно.

Страница \texttt{HistoryPage} отображает таблицу выполненных тестов с возможностью выбора одного теста для режима \texttt{per\_test} или от двух до пяти~--- для режима \texttt{series}. Страница \texttt{ComparisonPage} визуализирует ответ \texttt{POST /api/comparison/analyze}: резюме, таблицу метрик~--- в режиме \texttt{per\_test} отклонения относительно первого подключения в порядке \texttt{connection\_ids} прогона, в режиме \texttt{series}~--- относительно базового прогона,~--- диаграммы, статистические тесты, аналитический отчёт и при \texttt{series} блок \texttt{parameter\_impacts}; экспорт результата в JSON выполняется кнопкой Download на той же странице. Страница \texttt{DatabaseStatePage} предоставляет панель ручного сохранения и восстановления состояния подключённых баз данных и оценки объёма данных.

\subsubsection{Управление состоянием}

Состояние приложения хранится в едином store, реализованном через библиотеку Zustand. Файл \texttt{lib/store.ts} определяет интерфейс \texttt{AppState} и создаёт store через \texttt{create<AppState>}. Store включает следующие группы данных: навигация~--- текущая страница и состояние боковой панели; конфигурация теста~--- режим, сценарий, идентификатор набора сценария нагрузки или шаблона, признак использования индексов, число виртуальных пользователей, итерации и время прогрева; текущий тест; история тестов; данные реального времени~--- временные ряды для каждого подключения; словари имён и типов подключений; параметры сравнительного анализа.

Метод \texttt{addRealtimeData} обеспечивает добавление точек временного ряда с ограничением буфера.

Метод \texttt{.slice(-100)} ограничивает массив последними 100 точками для каждого подключения, что предотвращает чрезмерное потребление памяти при длительных тестах и обеспечивает плавную отрисовку графиков. Ключом в словаре \texttt{realtimeData} выступает уникальный идентификатор подключения, что позволяет отображать метрики для нескольких экземпляров одного типа СУБД независимо.

\subsubsection{Реализация WebSocket-хука}

Хук \texttt{useTestWebSocket} файла \texttt{hooks/use-test-websocket.ts} обеспечивает двустороннюю связь с сервером для получения метрик в реальном времени. При монтировании компонента с валидным \texttt{testId} автоматически устанавливается WebSocket-соединение по адресу \texttt{/ws/test/\{testId\}}. Каждые 30 секунд отправляется ping-сообщение для предотвращения разрыва соединения из-за таймаута.

При разрыве соединения во время выполнения теста производится автоматическое переподключение с интервалом 3 секунды. Переподключение выполняется только при статусе \texttt{running}, что предотвращает бесконечные попытки подключения после завершения теста. При получении сообщения типа \texttt{metrics} данные преобразуются в \texttt{TimeSeriesPoint} и добавляются в store через \texttt{addRealtimeData}, что вызывает реактивное обновление компонентов дашборда.

\subsubsection{Модульный API-клиент}

API-клиент реализован как класс \texttt{ApiClient} файла \texttt{lib/api/client.ts}. Базовый URL определяется переменной окружения \texttt{NEXT\_PUBLIC\_API\_URL}. Метод \texttt{request<T>} централизует обработку ошибок: при получении HTTP-статуса, отличного от 2xx, из ответа извлекается сообщение об ошибке и выбрасывается исключение \texttt{Error}. Все методы клиента типизированы через дженерики TypeScript, что обеспечивает проверку корректности данных на этапе компиляции.

Класс включает доменные методы для каждой группы эндпоинтов: методы управления подключениями и их профилями, методы работы с профилями схемы, методы управления группами баз данных, а также методы запуска тестов, получения истории и запуска сравнительного анализа. Экземпляр \texttt{apiClient} экспортируется как singleton; компоненты импортируют методы через barrel-экспорт \texttt{lib/api.ts}.

\subsubsection{Визуализация данных}

Визуализация метрик реализована с использованием библиотек Recharts и Plotly. Секундные ряды дашборда отображаются в переиспользуемом компоненте \texttt{TimeSeriesChart}; подписи пропускной способности зависят от \texttt{workload\_mode} прогона. На странице сравнительного анализа для диаграмм размаха и расширенных статистических графиков применяется Plotly.

Стилизация интерфейса выполнена через Tailwind CSS v4 с использованием дизайн-токенов для цветов, отступов и размеров шрифтов; поддерживается тёмная тема через \texttt{className="dark"} на корневом элементе. Компоненты интерфейса построены на базе Radix UI через библиотеку shadcn/ui. Уведомления пользователю отображаются через библиотеку Sonner в виде всплывающих сообщений.

\subsection{Контейнеризация и развёртывание}

Развёртывание системы организовано через Docker Compose. Основной файл \texttt{docker-compose.yml} определяет три сервиса и представлен в таблице~\ref{tab:docker_services}.

\begin{table}
    \centering
    \caption{Сервисы Docker Compose}
    \label{tab:docker_services}
    \begin{tabularx}{\textwidth}{|>{\RaggedRight}X|>{\RaggedRight}X|>{\RaggedRight}X|>{\RaggedRight}X|}
        \hline
        Сервис & Образ/источник & Назначение & Порт \\
        \hline
        project-db & postgres:15-alpine & База данных проекта & 5433:5432 \\
        \hline
        backend & Dockerfile на Python & FastAPI REST API и WebSocket-сервер & 8000:8000 \\
        \hline
        frontend & Dockerfile на Node.js & Next.js веб-интерфейс & 3000:3000 \\
        \hline
    \end{tabularx}
\end{table}

Сервис \texttt{project-db} развёртывается на базе образа \texttt{postgres:15-alpine}. Сервис настроен с именованным томом \texttt{project\_db\_data} для персистентности данных, healthcheck через \texttt{pg\_isready} с интервалом 5 секунд и политикой перезапуска \texttt{unless-stopped}. Порт~5433 используется для обеспечения совместимости с окружениями, где порт~5432 уже занят.

Сервис \texttt{backend} собирается из \texttt{backend/Dockerfile} и запускает FastAPI-приложение через uvicorn. Зависимость \texttt{depends\_on} с условием \texttt{service\_healthy} гарантирует, что бэкенд не стартует до полной готовности базы данных. Параметр \texttt{extra\_hosts: host.docker.internal:host-gateway} обеспечивает доступность хостовых СУБД из контейнера. \texttt{backend/Dockerfile} устанавливает системные пакеты \texttt{postgresql-client} и \texttt{mariadb-client}, что обеспечивает доступность CLI-утилит \texttt{pg\_dump}/\texttt{pg\_restore} и \texttt{mariadb-dump}/\texttt{mariadb} внутри контейнера для нативной стратегии управления состоянием баз данных.

Сервис \texttt{frontend} собирается из \texttt{frontend/Dockerfile} через многостадийную сборку: на первой стадии выполняется \texttt{next build} в standalone-режиме, на второй~--- копируется минимальный набор артефактов. Переменные \texttt{NEXT\_PUBLIC\_API\_URL} и \texttt{NEXT\_PUBLIC\_WS\_URL} задаются на этапе сборки и указывают на бэкенд по адресу \texttt{localhost:8000}. Все три сервиса объединены в пользовательскую сеть \texttt{loadtest-network} типа bridge.

Тестовые СУБД развёртываются отдельно через собственные Docker Compose-файлы в директориях \texttt{../sakila/} и \texttt{../pagila/}~--- соседних относительно корня репозитория. MySQL с схемой Sakila и PostgreSQL с схемой Pagila при первом запуске автоматически инициализируют схему. Оба контейнера доступны из бэкенда через хост \texttt{host.docker.internal}.

\subsubsection{Инициализация и миграции}

Процедура старта бэкенда включает миграции Alembic и \texttt{LogicalScenarioBootstrap}; она описана выше при рассмотрении менеджера контекста lifespan. Цепочка Alembic включает 17 версионируемых миграций, последовательно развивающих модель данных: от базовых таблиц прогонов, результатов и подключений~--- к индексам сценариев, профилям схем, группам баз данных и транзакционным наборам OLTP. Повторные запуски не создают дублирующих записей встроенных сущностей.

Конфигурация системы задаётся через файл \texttt{.env}: переменная \texttt{HISTORY\_DATABASE\_URL} задаёт подключение к базе данных проекта, \texttt{ENCRYPTION\_KEY}~--- ключ шифрования паролей подключений, \texttt{API\_HOST} и \texttt{API\_PORT}~--- настройки API-сервера, \texttt{NEXT\_PUBLIC\_WS\_URL} и \texttt{NEXT\_PUBLIC\_API\_URL}~--- адреса для фронтенда. Файл \texttt{.env} не включается в систему контроля версий; вместо него поставляется шаблон \texttt{env.example}.

Полный запуск системы сводится к одной команде: \texttt{docker compose up -d --build}. После её выполнения поднимаются три сервиса из таблицы~\ref{tab:docker_services}; при старте бэкенда автоматически выполняются шаги lifespan. Веб-интерфейс доступен по адресу \texttt{http://localhost:3000}, документация API~--- по адресу \texttt{http://localhost:8000/docs}.

Описанная программная реализация охватывает все компоненты, спроектированные в предыдущей главе, и обеспечивает полную функциональность системы нагрузочного тестирования.

\subsection{Выводы по главе}

В настоящей главе представлена программная реализация системы нагрузочного тестирования. Серверная часть разработана на Python~3.11 с использованием FastAPI, асинхронных драйверов СУБД и SQLAlchemy~2.0; управление схемой автоматизировано цепочкой из 17 миграций Alembic. Реализованы пять репозиториев слоя доступа к данным, модуль генерации нагрузки с поддержкой одиночных SQL-запросов и транзакционных последовательностей операций, слой диалектов СУБД для унифицированного сбора внутренних метрик, а также модуль сравнительного анализа с двумя режимами работы и полным набором статистических критериев. Клиентская часть на основе Next.js~16 / React~19 предоставляет восемь страниц интерфейса с визуализацией метрик в реальном времени через WebSocket, Recharts и Plotly. Развёртывание реализовано через Docker Compose с автоматической инициализацией и идемпотентным bootstrap. На основе реализованной системы в следующей главе проводится экспериментальная проверка и анализ результатов.
